
\chapter{Resultados y conclusiones}
\label{cha:resultados-conclusiones}

Este capítulo cierra el documento contrastando los resultados obtenidos con los
objetivos planteados en el capítulo introductorio, sintetizando las
aportaciones del trabajo y delimitando con precisión las limitaciones
detectadas durante la experimentación, así como las líneas naturales de
continuación.

\section{Síntesis de las aportaciones}

El Trabajo de Fin de Máster ha producido tres aportaciones articuladas en torno
a la extensión del corpus WebNLG al español. En primer lugar, una
\textbf{plataforma web propia}, \texttt{lanbench}, que materializa un flujo
\textit{Human-in-the-Loop} de anotación y revisión sobre un modelo de
\textit{datasets}, secciones, entradas y permisos. La arquitectura por capas
descrita en el capítulo~\ref{cha:experimentacion} se ha implementado y se ha
verificado mediante una batería de pruebas unitarias y de integración que
cubren la API HTTP, la persistencia con Prisma sobre MariaDB y la asistencia
del componente LLM.

En segundo lugar, una \textbf{integración con modelos de lenguaje de gran
	escala} en dos modalidades operativas, configurables por \textit{dataset}: una
modalidad \textit{corrector} que opera como asistente del anotador y del
revisor, y una modalidad \textit{creativo} que genera verbalizaciones
candidatas a partir de las tripletas RDF. El componente se apoya en un cliente
compatible con la API de OpenAI, lo que permite intercambiar proveedores (Groq,
Google AI Studio, Anthropic) sin intervención sobre el código de servicio. Las
claves se almacenan cifradas por \textit{dataset}, en línea con la historia
US-31, y no se devuelven al cliente en claro en ningún momento.

En tercer lugar, un \textbf{paquete experimental reproducible}: dos corpora de
50 entradas derivados de forma determinista de los ficheros maestro y de
\texttt{ru\_dev.xml}; dos \textit{scripts} de evaluación multi-proveedor con
\textit{back-off} sensible al proveedor y fusión incremental del JSON de
salida; ocho ficheros de prueba sin acceso a red que verifican la consistencia
de los corpora, su alineación
\textit{input}~$\leftrightarrow$~\textit{esperado} y la no exposición de claves
en las salidas; y la totalidad de los \textit{system prompts} empleados,
recogidos íntegramente en el anexo correspondiente.

\section{Cumplimiento de los objetivos específicos}
\label{sec:cierre-objetivos}

A continuación se contrasta cada objetivo específico enunciado en
~\ref{sec:proposito-objetivos-trabajo} con el grado de cumplimiento alcanzado.

\subsection*{OE1: \textit{pipeline} de corrección automática}

El objetivo de diseñar e implementar un \textit{pipeline} de corrección que
detecte y subsane los errores introducidos por la traducción se ha cubierto
\textbf{en su totalidad funcional y se ha cuantificado experimentalmente}. El
\textit{pipeline} \texttt{spanish-service.checkBatch} combina una etapa
determinista basada en reglas, una llamada al modelo de lenguaje y una
verificación final de cobertura. La calibración dirigida de las constantes
léxicas eleva el acuerdo con la referencia humana del 60\,\% al 85\,\% en el
corpus inicial de 20 entradas, y el comportamiento del \textit{pipeline} sobre
el corpus extendido de 50 entradas confirma el sesgo previsto hacia las
\textit{warnings} de fluidez. Las líneas de mejora identificadas en
§\ref{sec:limitaciones-trabajo-futuro} señalan los pasos concretos para
estrechar la brecha residual.

\subsection*{OE2: reproducibilidad del proceso}

El objetivo de reproducibilidad se considera \textbf{cumplido} en sus tres
dimensiones. La configuración de cada experimento se documenta en el cuerpo del
capítulo de experimentación, incluyendo proveedores, modelos, \textit{prompts}
y semillas. Los corpora se regeneran de forma determinista a partir de los
ficheros maestro; los \textit{scripts} de evaluación producen ficheros JSON
acumulables, de modo que pasadas parciales no sobrescriben las anteriores; y el
anexo de \textit{system prompts} reproduce literalmente las instrucciones de
sistema empleadas. El procedimiento operativo está descrito tanto en el cuerpo
del capítulo como en el anexo de reproducción.

\section{Conclusiones del experimento}

Los resultados de la experimentación, descritos en detalle en el
capítulo~\ref{cha:experimentacion}, permiten extraer cuatro conclusiones
principales que se enuncian a continuación de forma sintética para facilitar la
lectura cruzada con el capítulo.

\begin{enumerate}
	\item \textbf{Calibración determinista}. El ajuste dirigido de las constantes léxicas
	      en \texttt{domain/spanish/lexicon.js} resulta desproporcionadamente eficaz: una
	      intervención acotada (ampliación de \texttt{COMPLETE\_SENTENCE\_MARKERS},
	      incorporación de participios pasivos a los patrones de relación y tabla
	      \texttt{ENTITY\_ALIASES}) explica la mayor parte de la mejora del 60\,\% al
	      85\,\%, sin introducir falsos negativos sobre las clases de error genuinas.
	\item \textbf{Brecha residual sobre \textit{warnings} de fluidez}. El descenso del
	      acuerdo del 85\,\% al 68\,\% en el corpus extendido se explica por el peso de
	      las \textit{warnings} de fluidez (10/50), no por una degradación intrínseca del
	      \textit{pipeline}. La línea de mejora natural es la regla de supresión en el
	      \textit{alert-merger} descrita en \S\ref{sec:limitaciones-trabajo-futuro}.
	\item \textbf{Comparativa entre proveedores}. La submuestra efectiva de Gemini 2.5
	      Flash alcanza un acuerdo del 91\,\% sobre 22 entradas, pero la cuota gratuita
	      impide una comparación estadísticamente significativa con Llama 3.3 70B. La
	      hipótesis~H3 queda parcialmente sustentada y se mantiene como pregunta abierta
	      a la espera de cuota pagada.
	\item \textbf{Validación end-to-end del flujo de generación}. El
	      \texttt{auto-annotation-service} produce JSON sintácticamente válido,
	      semánticamente plausible y clasificado con severidad \texttt{ok} por el propio
	      validador del sistema en las pasadas en que el proveedor respondió. El bloqueo
	      cuantitativo es operativo (cuota), no funcional.
\end{enumerate}

\section{Limitaciones del alcance actual y trabajo futuro}
\label{sec:limitaciones-trabajo-futuro}

Las limitaciones detectadas se agrupan en tres familias, cada una de las cuales
acota una línea de trabajo futuro inmediata.

\subsection{Funcionalidades pendientes en la aplicación}

El documento de historias de usuario identifica capacidades que, en el estado
actual del repositorio, no se encuentran cubiertas y que se contemplan como
líneas de trabajo futuro:

\begin{itemize}
	\item \textbf{Configuración dinámica de criterios de evaluación} (US-24): se conserva
	      en la memoria por su interés funcional, pero su implementación queda como
	      trabajo pendiente. El alcance propuesto incluye la creación, edición y
	      versionado de criterios desde la página de administración.
	\item \textbf{Segunda ronda de revisión asociada al campo de \textit{revisiones
			      adicionales tras corrección}}: la opción se persiste con el \textit{dataset}
	      pero todavía no dispara ningún flujo adicional. El diseño previsto encadena
	      automáticamente una segunda revisión sobre las entradas corregidas, con
	      asignación a un revisor distinto del primero.
	\item \textbf{URL de invitación parametrizables por rol y por uso}: descritas en el
	      capítulo de la aplicación dentro del bloque de permisos y gobernanza, quedan
	      pendientes de implementación.
\end{itemize}

\subsection{Limitaciones de la evaluación experimental}

La experimentación reportada en el capítulo~\ref{cha:experimentacion} presenta
tres limitaciones operativas relevantes que conviene explicitar. La primera es
el \textbf{tamaño muestral} (50 entradas), que limita la significancia
estadística de la comparación entre proveedores: la ampliación a varios cientos
de entradas, distribuidas en sesiones que no saturen la cuota gratuita,
permitiría contrastar la hipótesis~H3 con un test de
McNemar~\cite{McNemar1947}. La segunda es la \textbf{convivencia con el flujo
	de generación} en la misma ventana de 24 horas, responsable del agotamiento de
la cuota diaria de Groq durante la pasada de
\texttt{eval-generation-quality.js}: bastaría con ejecutar las dos evaluaciones
en días separados o promover Groq al \textit{Dev Tier} para desbloquear esta
limitación. La tercera es la \textbf{cobertura por familia de modelos}: la
generalización a otras familias (Claude, Mistral, modelos locales) queda
pendiente.

\subsection{Mejoras del \textit{pipeline} sugeridas por el análisis de errores}

El análisis de los desacuerdos del capítulo~\ref{cha:experimentacion}
identifica tres líneas concretas de mejora del \textit{pipeline}:

\begin{enumerate}
	\item Ampliación de la tabla \texttt{ENTITY\_ALIASES} en
	      \texttt{domain/spanish/lexicon.js} con los pares pendientes (Seville/Sevilla,
	      Buenos\_Aires/Argentina, entre otros), responsables de los falsos positivos de
	      \texttt{missing\_triple} sobre entidades correctamente localizadas.
	\item Ampliación de \texttt{COMPLETE\_SENTENCE\_MARKERS} con verbos secundarios
	      (\textit{hacía}, \textit{trabajaba}, \textit{incorpora}, \textit{desempeña},
	      \textit{efectúa}, \textit{procede a}) que aparecen de forma recurrente en
	      verbalizaciones poco naturales.
	\item Introducción de una \textbf{regla de supresión} en el \textit{alert-merger}:
	      cuando el detector haya emitido ya una alerta de
	      \texttt{unnatural\_expression}, suprimir las alertas posteriores de
	      \texttt{incomplete\_sentence} o \texttt{semantic\_mismatch} sobre la misma
	      frase. Esta regla cancela el patrón dominante de falsos positivos observado
	      sobre el corpus de 50 entradas.
\end{enumerate}

\subsection{Plan de pruebas manuales de sistema con usuarios reales}
\label{sec:future-work-pruebas-usuario}

Las pruebas automatizadas descritas en el capítulo de pruebas validan la
corrección funcional de los flujos por separado, pero no su \textbf{aceptación}
por parte de los tres roles operativos del sistema ---administrador, anotador y
revisor--- en condiciones realistas de uso. Como complemento natural a esa
batería, se propone un protocolo de pruebas manuales de sistema con usuarios
reales, alineado con la clasificación clásica de validación
empírica~\cite{Wohlin2012} y diseñado para ejecutarse en una única iteración
acotada en tiempo. El plan se documenta a continuación con suficiente nivel de
detalle para que pueda ejecutarse sin diseño adicional.

\paragraph{Objetivos.} La campaña persigue tres objetivos diferenciados:
(i)~confirmar que cada rol puede completar los escenarios críticos sin
intervención del experimentador; (ii)~cuantificar la usabilidad percibida de la
plataforma mediante una escala estandarizada; y (iii)~recoger observaciones
cualitativas que orienten la siguiente iteración de la interfaz.

\paragraph{Participantes.} La muestra mínima objetivo se cifra en
\textbf{ocho participantes} distribuidos en los tres perfiles:

\begin{itemize}
	\item 1--2 participantes con \textbf{perfil de administrador}, con experiencia previa
	      en gestión de campañas de anotación o de localización.
	\item 3--4 participantes con \textbf{perfil de anotador} hispanohablantes, con o sin
	      experiencia previa en traducción.
	\item 2--3 participantes con \textbf{perfil de revisor} hispanohablantes con criterio
	      lingüístico explícito.
\end{itemize}

La selección procurará incorporar al menos dos variantes dialectales distintas
del español (peninsular, rioplatense, mexicana o caribeña), de manera coherente
con la preocupación dialectal expresada en
\S\ref{sec:future-work-elo-social}. La participación se regulará mediante
consentimiento informado por escrito y se remunerará conforme a los criterios
de remuneración justa discutidos
en~\S\ref{sec:etica-crowdsourcing}~\cite{Whiting2019_FairWork,Shmueli2021_BeyondFairPay},
con el tratamiento de datos personales conforme al RGPD~\cite{GDPR2016}.

\paragraph{Escenarios de prueba.} Los escenarios cubren los caminos críticos
asociados a cada rol y se enumeran en la
Tabla~\ref{tab:escenarios-usuarios-reales}. Cada escenario es ejecutable de
forma independiente sobre una instancia limpia de la plataforma desplegada con
\texttt{docker compose up} y poblada con el \textit{seed} de la base de datos.

\begin{table}[ht]
	\centering
	\small
	\begin{tabular}{llp{6.0cm}}
		\hline
		\textbf{ID} & \textbf{Rol}   & \textbf{Escenario}                                                                                  \\
		\hline
		S1          & Administrador  & Crear un \textit{dataset}, configurar proveedor LLM y asignar revisores                             \\
		S2          & Administrador  & Generar una URL de invitación, enviarla y validar el acceso resultante                              \\
		S3          & Administrador  & Exportar los resultados de un \textit{dataset} completado y verificar la integridad del fichero     \\
		S4          & Anotador       & Aceptar una invitación, completar diez verbalizaciones y retomar el trabajo tras desconectarse      \\
		S5          & Anotador       & Solicitar una verbalización asistida por LLM, editarla y enviarla                                   \\
		S6          & Revisor        & Revisar las verbalizaciones asignadas, asignar severidades y dejar comentarios                      \\
		S7          & Revisor        & Detectar una falsa aceptación del LLM (caso preparado) y reclasificarla                             \\
		S8          & Multi-rol      & Recorrido completo desde el alta del \textit{dataset} hasta la exportación tras revisión            \\
		\hline
	\end{tabular}
	\caption{Escenarios de prueba manual con usuarios reales, derivados de las
		historias de usuario consolidadas en el capítulo de la aplicación.}
	\label{tab:escenarios-usuarios-reales}
\end{table}

\paragraph{Protocolo de sesión.} Cada sesión se estructura en cuatro bloques y
está acotada a 75~minutos:

\begin{enumerate}
	\item \textbf{Briefing} (5~min): presentación del propósito, firma del
	      consentimiento y verificación del entorno.
	\item \textbf{Ejecución cronometrada con \textit{think-aloud}} (50~min): el
	      participante completa los escenarios asignados a su rol verbalizando sus
	      decisiones; el experimentador observa sin intervenir, salvo bloqueo absoluto.
	\item \textbf{Cuestionario post-sesión} (10~min): escala
	      \textbf{System Usability Scale} (10~ítems, valor agregado en el rango 0--100) e
	      ítems Likert específicos sobre claridad del flujo y confianza en los veredictos
	      del LLM.
	\item \textbf{Entrevista semiestructurada} (10~min): tres preguntas abiertas sobre
	      mayor dificultad encontrada, mayor sorpresa positiva y propuesta de mejora
	      prioritaria.
\end{enumerate}

\paragraph{Métricas e instrumentos.} Para cada participante y escenario se
registran:

\begin{itemize}
	\item \textbf{Tasa de éxito} del escenario (binaria, con o sin asistencia).
	\item \textbf{Tiempo en tarea} (segundos) desde el inicio hasta la finalización.
	\item \textbf{Número de errores del usuario} (acciones equivocadas que requieren
	      \textit{undo} o reinicio del flujo).
	\item \textbf{Número de solicitudes de ayuda} dirigidas al experimentador.
	\item \textbf{Puntuación SUS} agregada por participante.
	\item \textbf{Observaciones cualitativas} extraídas del \textit{think-aloud}, codificadas
	      en categorías \textit{a posteriori}.
\end{itemize}

\paragraph{Criterios de aceptación.} Se considera que un escenario supera la
prueba cuando se cumplen, sobre el conjunto de participantes asignados:

\begin{itemize}
	\item Tasa de éxito sin asistencia $\geq$ 90\,\% en los escenarios marcados como
	      críticos (S1, S4, S6, S8) y $\geq$ 75\,\% en los restantes.
	\item Puntuación SUS media $\geq$ 70 por rol, umbral correspondiente a una
	      usabilidad calificada como \emph{buena} en la literatura especializada.
	\item Ausencia de defectos bloqueantes (severidad crítica) en el seguimiento de
	      incidencias durante la campaña.
\end{itemize}

\paragraph{Tratamiento de datos y entregables.} Las sesiones se grabarán
---previo consentimiento--- y los datos brutos se anonimizarán antes de su
análisis. El entregable de la campaña consta de un informe que reporta las
métricas agregadas, el listado priorizado de defectos detectados y un plan de
mitigación por iteración, incorporado como nuevo anexo al presente trabajo o
como adenda en una versión posterior. El protocolo, los formularios y la
plantilla de informe se publicarán junto al código fuente para favorecer su
reutilización en posteriores rondas de validación.

\subsection{Líneas de continuación a medio plazo}

Más allá de las acciones inmediatas, el trabajo abre tres líneas naturales de
continuación. La primera es la \textbf{producción del corpus a escala plena}
mediante una campaña de \textit{crowdsourcing} controlada, en línea con las
consideraciones éticas y de remuneración descritas en
\S\ref{sec:etica-crowdsourcing}. La segunda es la \textbf{evaluación cruzada
	con la rúbrica MQM} en lugar de la actual rúbrica ternaria, lo que habilitaría
la comparabilidad con la literatura especializada y la agregación ponderada de
severidades. La tercera es la \textbf{publicación formal de la extensión al
	español del \textit{benchmark}} en una plataforma como Hugging Face o
GERBIL/BENG, lo que requiere consolidar el formato extendido y acompañarlo de
la \textit{data card} prevista por la directiva GEM.

\subsection{Evaluación social en bucle: una metaheurística para anotadores
	y modelos generativos}
\label{sec:future-work-elo-social}

Las tres líneas anteriores asumen el esquema clásico de evaluación: rúbrica
prefijada, anotadores con asignación administrativa y agregación por
mayoría dentro de cada entrada. Este esquema reproduce, sin corregirlo, el
error individual de cada revisor y no genera una métrica de calidad por
\emph{anotador} o por \emph{modelo}: la calidad sólo se observa en
agregado, a posteriori y con coste alto. Una línea de continuación más
ambiciosa, complementaria a las anteriores, consiste en sustituir esa
agregación estática por una metaheurística social que aprenda de forma
continua tanto la calidad de cada revisor humano como la de cada modelo
generativo, en el espíritu de los esquemas de \emph{community-based
	moderation} popularizados por las \textit{Community Notes} de X
~\cite{Wojcik2022} y de los \textit{leaderboards} dinámicos de modelos
basados en preferencia humana~\cite{Chiang2024_ChatbotArena}.

\paragraph{Fundamento.} La intuición se apoya en tres bloques bien
establecidos. El primero es la \textbf{sabiduría de las multitudes}
~\cite{Surowiecki2004_Wisdom}: cuando los juicios son diversos,
independientes y se agregan mediante un mecanismo adecuado, el ruido
individual tiende a cancelarse y el agregado supera al juicio mediano
individual. El segundo es el conjunto de \textbf{sistemas de
	\textit{rating} por comparación pareada} ---Elo en
ajedrez~\cite{Elo1978_Rating}, Bradley--Terry~\cite{BradleyTerry1952} y
TrueSkill~\cite{Herbrich2007_TrueSkill}---, que infieren una puntuación
latente de habilidad a partir de victorias y derrotas sobre rivales,
ponderando la información por la incertidumbre acumulada. El tercero es
el \textbf{\textit{bridging-based ranking}} de las
\textit{Community Notes}~\cite{Wojcik2022}, que descompone las
valoraciones en dos factores latentes ---utilidad y polarización--- y premia
exclusivamente las contribuciones cuya utilidad percibida atraviesa
poblaciones de votantes con sesgos opuestos. A estos tres bloques se
añade, en el plano del anotador individual, la \textbf{estimación
	bayesiana de la competencia} del anotador frente al consenso latente, en
la línea de MACE~\cite{Hovy2013_MACE}, que permite ponderar votos por
fiabilidad estimada y detectar comportamientos de \textit{spam} o
adhesión sistemática.

\paragraph{Por qué la mayoría es un juez legítimo en lengua, pero no en
	ciencia.} Una objeción inmediata es que la mayoría no es, en general,
un buen árbitro de la verdad: en ciencia, la corrección de una
afirmación es independiente del número de personas que la suscriban. El
caso del idioma, sin embargo, presenta una asimetría epistemológica
relevante. Una verbalización en español no es ``correcta'' por
adecuación a una realidad externa, sino \emph{útil} en la medida en que
una comunidad de hablantes la entiende y la acepta como natural. Esta es
la posición descriptivista que ha guiado a la lingüística académica
desde Saussure y que se contrapone explícitamente al prescriptivismo
institucional ---de la RAE al manual de estilo--- sin negar su valor
operativo en la enseñanza y la edición. En este marco, una mayoría
amplia de hablantes diversos que coincide en que una frase ``se
entiende'' o ``no se entiende'' \emph{es} la señal que define el
fenómeno, no una aproximación a una señal externa. La aplicabilidad de
la sabiduría de las multitudes al juicio lingüístico es, por tanto,
sustantivamente más sólida que a juicios fácticos comparables.

\paragraph{Diseño propuesto sobre \texttt{lanbench}.} La materialización
de esta línea en una versión sucesiva de la plataforma admite la
siguiente arquitectura mínima:

\begin{enumerate}
	\item \textbf{Producción de pares}. Para cada entrada del
	      \textit{dataset}, varios productores ---humanos y modelos creativos
	      configurados con \emph{system prompts} distintos--- generan
	      verbalizaciones candidatas. La plataforma ya soporta el almacenamiento
	      por anotador (US-23) y la asignación de modelo por
	      \textit{dataset}; el cambio consiste en levantar la restricción de una
	      sola verbalización canónica por entrada durante esta fase
	      exploratoria.
	\item \textbf{Votación pareada de revisores}. La pantalla de revisión
	      sustituye la rúbrica ternaria por una comparación A--B entre
	      verbalizaciones de la misma entrada, sin información sobre el
	      productor. El revisor escoge la traducción que ``entiende mejor''
	      o ``percibe como más natural'' y, opcionalmente, marca un empate.
	      Esta interfaz reduce la carga cognitiva respecto a la rúbrica
	      explícita y ha demostrado ser la modalidad más fiable para
	      capturar preferencia humana en evaluación
	      \textit{open-ended}~\cite{Chiang2024_ChatbotArena}.
	\item \textbf{Agregación con \textit{rating} dinámico}. Cada
	      comparación actualiza un \textit{rating} por productor ---humano o
	      modelo--- mediante Bradley--Terry o TrueSkill, con la incertidumbre
	      explícita por número de combates. El \textit{leaderboard}
	      resultante es la \textbf{métrica de calidad en tiempo real} de
	      cada agente generativo, y elimina la necesidad de campañas
	      \textit{ad hoc} de evaluación de modelos cada vez que aparece un
	      nuevo proveedor.
	\item \textbf{Filtro \textit{bridging-based} contra polarización
		      dialectal}. La aplicación directa de Elo penaliza variantes
	      minoritarias del español ibérico, rioplatense o caribeño, porque
	      las votaciones se acumulan donde está la mayoría de anotadores.
	      Aplicar la descomposición en factores latentes utilidad/polarización de
	      las \textit{Community Notes}~\cite{Wojcik2022} permite identificar
	      qué verbalizaciones convencen \emph{a través} de variantes
	      regionales ---y por tanto generalizan--- y separarlas de las que
	      simplemente reflejan el dialecto mayoritario del panel.
	\item \textbf{Competencia del revisor y peso del voto}. En paralelo se
	      mantiene una estimación bayesiana de la competencia de cada
	      revisor frente al consenso emergente, al estilo de
	      MACE~\cite{Hovy2013_MACE}. El peso de su voto en las
	      actualizaciones futuras escala con esta competencia, lo que abre
	      la puerta a un \textbf{esquema de incentivos económicos
		      desigualitario} ---remuneración proporcional al \textit{rating} del
	      anotador--- alineado con la discusión ética y de remuneración
	      justa de
	      \S\ref{sec:etica-crowdsourcing}~\cite{Shmueli2021_BeyondFairPay,Whiting2019_FairWork}.
\end{enumerate}

\paragraph{Extensión al \textit{pipeline} de corrección.} La aplicación
del mismo esquema a la revisión ---es decir, a la identificación de
errores y no a su \emph{generación}--- es conceptualmente más exigente. La
calidad de una corrección no se reduce a una preferencia binaria entre
dos versiones, sino que involucra la detección de errores categóricos
(omisión, inserción, alucinación, agramaticalidad). Una posible
adaptación consiste en convertir cada divergencia entre revisores en una
votación pareada sobre el \emph{diagnóstico} ---``¿qué revisor ha
clasificado mejor este caso?''--- evaluada por un panel ciego, y agregar
los \textit{ratings} de detección por separado de los de generación. El
\textit{pipeline} automatizado del capítulo~\ref{cha:experimentacion}
encajaría como un revisor más dentro de este panel, lo que daría una
medida directa, dinámica y comparable de su rendimiento frente al
humano sin necesidad de una nueva campaña de etiquetado de oro.

\paragraph{Riesgos y mitigaciones reconocidas.} Esta línea hereda los
riesgos conocidos de los esquemas de \textit{rating} crowdsourced:
inicialización fría con varianza alta hasta acumular combates
suficientes, vulnerabilidad a coaliciones coordinadas, y el ya
mencionado sesgo hacia la variante dialectal mayoritaria. La
literatura ofrece mitigaciones específicas para cada uno ---incertidumbre
explícita de TrueSkill~\cite{Herbrich2007_TrueSkill}, detección de
manipulación inherente al \textit{bridging}~\cite{Wojcik2022}, refuerzo
explícito del muestreo de revisores en variantes minoritarias---, pero su
calibración para un corpus en español queda abierta como objeto de
investigación. El interés principal de esta línea no es por tanto su
inmediatez, sino que convierte la propia plataforma \texttt{lanbench}
en un \textit{benchmark} vivo de modelos generativos en español, donde
la calidad de cada proveedor se mide de forma continua y comparable, y
donde la calidad del trabajo humano de anotación se hace, por primera
vez en este flujo, observable y remunerable de forma diferencial.

\section{Reflexión final}

El trabajo no pretende sustituir el juicio humano sobre la calidad lingüística,
sino integrar el modelo de lenguaje como asistente sistemático en el bucle de
anotación, dejando la decisión última al revisor. La experimentación confirma
que el balance entre coste material y coste laboral favorece holgadamente esta
estrategia: el coste \textit{token} de una pasada completa sobre un
\textit{dataset} de tamaño WebNLG-2020 es inferior a un euro en ambos
proveedores, mientras que el coste laboral humano lo excede en más de dos
órdenes de magnitud. En este balance se sostiene la conveniencia operativa del
paradigma \textit{Human-in-the-Loop} y, con él, la utilidad de plataformas como
\texttt{lanbench} para soportar la producción de corpus lingüísticos en lenguas
con representación insuficiente en los \textit{benchmarks} de referencia.
